% Axial form factor comparison (LQCD vs deuterium): figures + chi2 tables in one
% file, each comparison's 6-panel figure followed by its two-row chi2/ndf table.
% Per-panel PDFs are co-located as fa_lfg_ha2018_{a..f}.pdf / fa_sf_incl_{a..f}.pdf;
% prepend your paper's path (e.g. plots/) to the \includegraphics keys.
% Requires: amsmath (for \text), graphicx, overpic, booktabs.
% Panels: (a) dpT (dalphaT<45 deg), (b) dpT (135<dalphaT<180), (c) dpT,
%         (d) dalphaT, (e) p_mu, (f) p_p.

% ======================= LFG + hA2018 : F_A LQCD vs deuterium =======================
\begin{figure*}[htbp]
\centering
\begin{minipage}{0.3\textwidth}
\begin{overpic}[width=\linewidth]{fa_lfg_ha2018_a.pdf}
\put(24,56){\footnotesize\text{(a)}}
\end{overpic}
\end{minipage}
%\hfill
\begin{minipage}{0.3\textwidth}
\begin{overpic}[width=\linewidth]{fa_lfg_ha2018_b.pdf}
\put(24,56){\footnotesize\text{(b)}}
\end{overpic}
\end{minipage}
%\hfill
\begin{minipage}{0.3\textwidth}
\begin{overpic}[width=\linewidth]{fa_lfg_ha2018_c.pdf}
\put(24,56){\footnotesize\text{(c)}}
\end{overpic}
\end{minipage}

%\hfill
\vspace{0.5cm}
\begin{minipage}{0.3\textwidth}
\begin{overpic}[width=\linewidth]{fa_lfg_ha2018_d.pdf}
\put(24,56){\footnotesize\text{(d)}}
\end{overpic}
\end{minipage}
\begin{minipage}{0.3\textwidth}
\begin{overpic}[width=\linewidth]{fa_lfg_ha2018_e.pdf}
\put(24,56){\footnotesize\text{(e)}}
\end{overpic}
\end{minipage}
%\hfill
\begin{minipage}{0.3\textwidth}
\begin{overpic}[width=\linewidth]{fa_lfg_ha2018_f.pdf}
\put(24,56){\footnotesize\text{(f)}}
\end{overpic}
\end{minipage}
\caption{Effect of the nucleon axial form factor on the MicroBooNE CC1$\mu$1p / CC1$\mu$Np argon cross sections with the N/LFG ground state and hA2018 FSI held fixed: $F_A^{\rm LQCD}$ (blue) versus $F_A^{\rm Deu}$ (orange), compared with MicroBooNE data (black points). Panels: (a) $\delta p_T$ for $\delta\alpha_T<45^\circ$ and (b) for $135^\circ<\delta\alpha_T<180^\circ$ (two slices of the 2D $\delta p_T$--$\delta\alpha_T$ measurement); (c) $\delta p_T$; (d) $\delta\alpha_T$; (e) $p_\mu$; (f) $p_p$. The LQCD form factor raises the predicted cross section and lowers $\chi^2/{\rm ndf}$ in every observable (see Table~\ref{tab:chisq_fa_lfg_ha2018}), favouring $F_A^{\rm LQCD}$.}
\label{fig:fa_lfg_ha2018}
\end{figure*}

\begin{table}[htbp]
\caption{$\chi^2/{\rm ndf}$ ($p$-value) relative to MicroBooNE data for Fig.~\ref{fig:fa_lfg_ha2018} (N/LFG ground state and hA2018 FSI held fixed).}
\label{tab:chisq_fa_lfg_ha2018}
\centering
\resizebox{\columnwidth}{!}{
\begin{tabular}{lrrrrr}
\toprule
Configuration & $\delta p_T$ vs. $\delta \alpha_T$ & $\delta p_T$ & $\delta \alpha_T$ & $p_{\mu}$ & $p_p$ \\
\midrule
N/LFG + $F_A^{\rm LQCD}$ + hA2018 & 35.98/49  (0.92) & 6.78/13  (0.91) & 3.47/7  (0.84) & 16.29/26  (0.93) & 16.05/15  (0.38)\\
N/LFG + $F_A^{\rm Deu}$ + hA2018  & 44.21/49  (0.67) & 9.59/13  (0.73) & 10.19/7  (0.18) & 20.89/26  (0.75) & 19.95/15  (0.17)\\
\bottomrule
\end{tabular}}
\end{table}

% ======================= SF + INCL : F_A LQCD vs deuterium =======================
\begin{figure*}[htbp]
\centering
\begin{minipage}{0.3\textwidth}
\begin{overpic}[width=\linewidth]{fa_sf_incl_a.pdf}
\put(24,56){\footnotesize\text{(a)}}
\end{overpic}
\end{minipage}
%\hfill
\begin{minipage}{0.3\textwidth}
\begin{overpic}[width=\linewidth]{fa_sf_incl_b.pdf}
\put(24,56){\footnotesize\text{(b)}}
\end{overpic}
\end{minipage}
%\hfill
\begin{minipage}{0.3\textwidth}
\begin{overpic}[width=\linewidth]{fa_sf_incl_c.pdf}
\put(24,56){\footnotesize\text{(c)}}
\end{overpic}
\end{minipage}

%\hfill
\vspace{0.5cm}
\begin{minipage}{0.3\textwidth}
\begin{overpic}[width=\linewidth]{fa_sf_incl_d.pdf}
\put(24,56){\footnotesize\text{(d)}}
\end{overpic}
\end{minipage}
\begin{minipage}{0.3\textwidth}
\begin{overpic}[width=\linewidth]{fa_sf_incl_e.pdf}
\put(24,56){\footnotesize\text{(e)}}
\end{overpic}
\end{minipage}
%\hfill
\begin{minipage}{0.3\textwidth}
\begin{overpic}[width=\linewidth]{fa_sf_incl_f.pdf}
\put(24,56){\footnotesize\text{(f)}}
\end{overpic}
\end{minipage}
\caption{As in Fig.~\ref{fig:fa_lfg_ha2018} but with the spectral-function ground state and INCL FSI held fixed: $F_A^{\rm LQCD}$ (blue) versus $F_A^{\rm Deu}$ (orange). The LQCD form factor again improves every observable (see Table~\ref{tab:chisq_fa_sf_incl}), but with INCL FSI both predictions remain well below the data, so varying $F_A$ alone does not reconcile the SF+INCL configuration with MicroBooNE.}
\label{fig:fa_sf_incl}
\end{figure*}

\begin{table}[htbp]
\caption{$\chi^2/{\rm ndf}$ ($p$-value) relative to MicroBooNE data for Fig.~\ref{fig:fa_sf_incl} (spectral-function ground state and INCL FSI held fixed).}
\label{tab:chisq_fa_sf_incl}
\centering
\resizebox{\columnwidth}{!}{
\begin{tabular}{lrrrrr}
\toprule
Configuration & $\delta p_T$ vs. $\delta \alpha_T$ & $\delta p_T$ & $\delta \alpha_T$ & $p_{\mu}$ & $p_p$ \\
\midrule
SF + $F_A^{\rm LQCD}$ + INCL & 66.61/49  (0.05) & 15.72/13  (0.26) & 5.47/7  (0.60) & 23.91/26  (0.58) & 32.25/15  (0.01)\\
SF + $F_A^{\rm Deu}$ + INCL  & 81.91/49  (0.00) & 23.94/13  (0.03) & 10.01/7  (0.19) & 35.08/26  (0.11) & 38.83/15  (0.00)\\
\bottomrule
\end{tabular}}
\end{table}
